Fix $n\ge1$, $N=3^n$, $\omega=e^{2\pi i/N}$; for $j\in\Z/N$ let $\lambda_j=\log|\sigma^j\eta_n|$ and $\hat\lambda_k=\sum_{j\in\Z/N}\lambda_j\,\omega^{-jk}$ ($k\in\Z/N$).

\emph{Step 1: $\hat\lambda_k=0$ for $3\mid k$.} By Lemma~\ref{lem:normone}(iii), $\lambda_j+\lambda_{j+N/3}+\lambda_{j+2N/3}=\log|\sigma^j\operatorname{Nr}_{B_{3,n}/B_{3,n-1}}\eta_n|=0$ for every $j$ ($\tau=\sigma^{N/3}$). For $k=3k'$ one has $\omega^{-(j+N/3)k}=\omega^{-jk}$, so grouping $\Z/N$ into the $N/3$ cosets of $\{0,N/3,2N/3\}$ gives $\hat\lambda_k=\sum_{j<N/3}\omega^{-jk}(\lambda_j+\lambda_{j+N/3}+\lambda_{j+2N/3})=0$ (\path{WeberP3Rel.dft_vanish_of_relnorm}).

\emph{Step 2: $\hat\lambda_k\ne0$ for $3\nmid k$.} Since $|\sin(2\pi m/q)|=\frac12|1-\zeta_q^{2m}|$, $\lambda_j=g(2a\cdot4^j)-g(2\cdot4^j)$ with $g(m)=\log|1-\zeta_q^m|$ on $(\Z/q)^\times=\{\pm1\}\times\langle4\rangle$ ($q=3N$, $a=1+N$, and $a\equiv4^{N/3}\pmod q$ by Lemma~\ref{lem:normone}(ii), \path{WeberP3Rel.four_pow_modEq}; the arguments of $g$, and below of $\chi_k$, are residue classes modulo $q$). As $g(-m)=g(m)$ and $-2\in\langle4\rangle$, write $-2\equiv4^{j_0}\pmod q$; then $\lambda_j=\Gamma(j+j_0+N/3)-\Gamma(j+j_0)$ with $\Gamma(j)=g(4^j)$, hence $\hat\lambda_k=\omega^{j_0k}(\omega^{Nk/3}-1)\hat\Gamma_k$ (shift rule of the twisted sum on $\Z/N$: \path{WeberP3Rel.twisted_sum_shift}, \path{WeberP3Rel.twisted_sum_shift_diff}) and $\omega^{Nk/3}=e^{2\pi ik/3}\ne1$ (\path{WeberP3Rel.pow_third_eq_one_iff}); so $\hat\lambda_k\ne0$ as soon as $\hat\Gamma_k\ne0$ (\path{WeberP3Rel.twisted_sum_shift_diff_ne_zero}). Let $\chi_k$ be the character of $(\Z/q)^\times=\{\pm1\}\times\langle4\rangle$ trivial on $-1$ with $\chi_k(4^j)=\omega^{-jk}$ (well defined since $4$ has order $N$ in $(\Z/q)^\times$ and $\omega^N=1$); it is even. \emph{Conductor.} Every proper divisor of $q=3^{n+1}$ divides $3^n$, so $\chi_k$ is imprimitive iff it factors through the reduction $(\Z/q)^\times\to(\Z/3^n)^\times$, i.e.\ iff it is trivial on the kernel of that reduction. The kernel is $\{t\in(\Z/q)^\times:t\equiv1\ (3^n)\}=\{1,\,1+3^n,\,1+2\cdot3^n\}=\{1,a,a^2\}$ (order $3$; $a=1+N$ has exact order $3$ by Lemma~\ref{lem:normone}(i), and $a^2\equiv1+2N$ as $N^2\equiv0\pmod q$). Since $a\equiv4^{N/3}\pmod q$ (Lemma~\ref{lem:normone}(ii), \path{WeberP3Rel.four_pow_modEq}) and $\chi_k$ is a function of the class modulo $q$, $\chi_k(a)=\omega^{-Nk/3}=e^{-2\pi ik/3}\ne1$ for $3\nmid k$. So $\chi_k$ is not trivial on the kernel: it is primitive of conductor exactly $q=3^{n+1}$. \emph{Conventions.} Throughout, $\omega=e^{2\pi i/N}$ and $\hat\lambda_k=\sum_j\lambda_j\omega^{-jk}$ (minus sign in the exponent; the same sign is used in the certificate script and in the read-only verifier, and the Lean statements of \path{WeberP3Rel} are proved for an arbitrary $\omega$ with $\omega^N=1$, so they cover either sign); only $|\hat\Gamma_k|$ enters below, through the absolute-value form of Washington's formula (Lemma~\ref{lem:D3}), so no Gauss-sum normalisation is used. Then $\hat\Gamma_k=\sum_{m\in\langle4\rangle}\chi_k(m)g(m)=\frac12\sum_{m\in(\Z/q)^\times}\chi_k(m)\log|1-\zeta_q^m|$, whose absolute value is $\frac{\sqrt q}2|L(1,\chi_k)|$ by Lemma~\ref{lem:D3} (Washington's formula, \cite{Wash}, Chapter~4), and $L(1,\chi_k)\ne0$ by Dirichlet's theorem. So $\hat\lambda_k\ne0$.

\emph{Step 3: rank $c$ for every $n$ and $1\le r\le n$.} Put $h=3^{n-r}$, $c=2\cdot3^{r-1}=\varphi(3^r)$ and $\xi=\omega^{h}$, a primitive $3^r$-th root of unity. For $a_0,\dots,a_{c-1}\in\R$ the vector $v=\sum_{i<c}a_iH(\sigma^{hi}\eta_n)$ has coordinates $v_j=\sum_ia_i\lambda_{j+hi}$ and transform $\hat v_k=A(\omega^{hk})\hat\lambda_k$, $A(x)=\sum_ia_ix^i$. If $v=0$ then $A(\xi^{k})=0$ for all $3\nmid k$ by Step~2; as $k$ runs over $(\Z/3^r)^\times$ the $\xi^k$ are the $c$ primitive $3^r$-th roots of unity, and $\deg A<c$ forces $A=0$ (\path{WeberP3Rel.eq_zero_of_vanish_on_primitiveRoots}, \path{WeberP3Rel.totient_three_pow}). Hence the $c$ vectors $H(\sigma^{hi}\eta_n)$ are linearly independent: (Rank) holds and $D_r^{(n)}>0$.

\emph{Step 4: exact covolume.} With Parseval's identity $\langle u,v\rangle=N^{-1}\sum_k\hat u_k\overline{\hat v_k}$ the Gram matrix of the $c$ vectors is
\[ G_{ii'}=\frac1N\sum_{k}|\hat\lambda_k|^2\,\omega^{hk(i-i')}=\sum_{b\in(\Z/3^r)^\times}W_{n,r}(b)\,\xi^{b(i-i')},\] where $W_{n,r}(b)=\frac1N\sum_{k\equiv b\ (3^r)}|\hat\lambda_k|^2>0$,
the terms $3\mid k$ being $0$ by Step~1 and the others positive by Step~2. \emph{Normalisation.} The left side $G_{ii'}=\sum_{j\in\Z/N}\lambda_{j+hi}\lambda_{j+hi'}$ is the Gram matrix of the log vectors themselves (the Gram route of the certificate computes it directly from the $\lambda_j$), and the right side is Parseval's identity for the unnormalised transform $\hat u_k=\sum_j u_j\omega^{-jk}$, whose inverse carries the factor $1/N$; this is the only place where $1/N$ enters, and it is the $1/N$ in the definition of $W_{n,r}(b)$ (the character-product route computes the right side from the $\hat\lambda_k$). Thus $G=V\operatorname{diag}(W)V^{*}$ with $V=(\xi^{bi})_{i<c,\ b\in(\Z/3^r)^\times}$, and $|\det V|^2=\prod_{b\ne b'}|\xi^{b}-\xi^{b'}|=|\operatorname{disc}\Phi_{3^r}|=|\operatorname{disc}\Q(\zeta_{3^r})|=3^{3^{r-1}(2r-1)}$ (Vandermonde over the $c$ primitive $3^r$-th roots of unity; the discriminant is Lemma~\ref{lem:disc}). Therefore $(D_r^{(n)})^2=\det G=3^{3^{r-1}(2r-1)}\prod_{b}W_{n,r}(b)$, which is the second route of the certificate \artPcert{}: the two routes are independent linear-algebraic evaluations from the same log profile $(\lambda_j)$ (Gram determinant; character product), and they agree to the ball radius in all $15$ rows.
